\chapter[Herramientas y tecnologías empleadas]{Herramientas y tecnologías\\empleadas}
\label{cap:herramientas}
Este capítulo presenta las herramientas y tecnologías que utilicé para construir
Neuma, verificar sus componentes, organizar el proyecto y preparar el informe.
La selección se apoyó en las necesidades de la aplicación y en propuestas técnicas
de los agentes que acepté y revisé durante el desarrollo. Mi experiencia previa
en aplicaciones Android era limitada y se concentraba en React Native; por ello,
la revisión del comportamiento obtenido tuvo un papel central en la aceptación
de las propuestas. El proceso de trabajo y supervisión se describe en el
capítulo~\ref{cap:metodologia}.

\section{Plataforma Android e interfaz}
\label{sec:herramientas-android}
Elegí Android como plataforma exclusiva y adopté Kotlin y Jetpack Compose para
el desarrollo nativo. Buscaba acceso directo a las capacidades de audio e
interacción del sistema y un entorno que facilitara el control de latencia y
la implementación de los juegos. También consideré una eventual conexión con
un teclado MIDI. Estos criterios orientaron la elección frente a continuar con
React Native, que conocía de experiencias previas. La decisión expresa una
expectativa técnica de integración; no deriva de una comparación experimental
de rendimiento entre ambos enfoques. La conexión MIDI permanece como posibilidad
considerada, fuera de la implementación descrita.

Jetpack Compose proporciona un enfoque declarativo para construir interfaces
Android \parencite{androidCompose}. En Neuma lo utilizo para representar pantallas
y estados de interacción, con componentes Material 3 adaptados a la dirección
artística del proyecto. La aplicación incorpora ViewModel y componentes Lifecycle
para vincular el estado de las actividades con la interfaz. Las corrutinas de
Kotlin permiten organizar operaciones asíncronas \parencite{kotlinCorrutinas};
en el código se utilizan, entre
otras tareas, para captura de audio y coordinación con los modelos de presentación.
Esta combinación sostiene una interfaz nativa y permite mantener las reglas
musicales y de adaptación en componentes separados.

La configuración fija un nivel mínimo de API 26, compilación y destino API 36,
y bytecode compatible con JVM 17. Para las bibliotecas Compose utiliza una
\textit{Bill of Materials} (BOM), que agrupa versiones compatibles entre sí
\parencite{androidComposeBom}. Los valores del proyecto se reúnen en la
tabla~\ref{tab:versiones-herramientas}; describen la configuración examinada,
sin implicar que todas las herramientas mantuvieran esa versión durante el trabajo.

\section{Audio, notación y almacenamiento local}
\label{sec:herramientas-audio-datos}
El acceso al micrófono utiliza \texttt{AudioRecord}, que permite leer muestras
desde la entrada de audio de Android. Para reproducir referencias y sonidos,
la aplicación utiliza \texttt{AudioTrack}, que admite el envío de búferes PCM
a la salida de audio \parencite{androidAudioRecord,androidAudioTrack}.
Estas APIs proporcionan acceso a las muestras que procesa el prototipo y a las
señales que genera o reproduce. Su incorporación mantiene el procesamiento
musical en el dispositivo; la latencia y la precisión deben analizarse bajo
las condiciones de prueba descritas en el capítulo~\ref{cap:pruebas}.

Sobre esas muestras se ejecuta una implementación propia de YIN en Kotlin.
El ajuste de dificultad también utiliza código propio, con un motor de inferencia
Mamdani y configuraciones por actividad. Son componentes del sistema cuya
selección e integración se apoyan en el marco conceptual del
capítulo~\ref{cap:marco}. El cálculo difuso se desarrolla en la
sección~\ref{sec:diseno-dda}, y el procesamiento de audio y la estimación de tono,
en la sección~\ref{sec:implementacion-audio}.

Para la notación de Rhythm Tap incorporé VexFlow, una biblioteca de grabado musical
que puede generar representaciones SVG y Canvas \parencite{vexflowDocumentacion}.
La integración utiliza SVG en una WebView con recursos empaquetados localmente,
incluidas las fuentes Bravura y Academico. La biblioteca permite generar la
partitura desde datos estructurados y resolver su disposición musical. La
WebView cumple esa función dentro de una aplicación nativa; la entrada táctil
de ejecución, el audio y la evaluación permanecen en los componentes Android.
El inventario de recursos conserva las versiones, procedencia y licencias de
la biblioteca y las fuentes.

Los resultados, el progreso y las configuraciones se almacenan en archivos JSON
del espacio privado de la aplicación. Para convertir las estructuras de Kotlin
a ese formato y recuperarlas se utiliza \texttt{kotlinx.serialization}
\parencite{kotlinSerializacion}.
El almacenamiento específico de la aplicación proporciona un ámbito interno
para estos archivos \parencite{androidArchivosPrivados}. En el prototipo, los
archivos respaldan la persistencia local y la exportación de registros en JSON
y ZIP. Esta solución respalda la conservación y consulta local de los datos
del prototipo.

\section{Construcción, pruebas y automatización}
\label{sec:herramientas-verificacion}
La construcción utiliza Gradle y Android Gradle Plugin. El repositorio incluye
Gradle Wrapper, que fija la distribución empleada y permite invocar la construcción
con una versión definida \parencite{gradleWrapper}. El SDK de Android aporta las
herramientas de plataforma y compilación, y Android Debug Bridge (ADB) interviene
en la comunicación con dispositivos o emuladores. Los scripts del proyecto
configuran el entorno de Java y Android y generan los paquetes APK de prueba.

Para las pruebas locales se utilizan JUnit y las utilidades de prueba de corrutinas.
Estas pruebas se ejecutan en la JVM de la computadora de desarrollo, de acuerdo
con la distinción de Android entre pruebas locales y pruebas que requieren el
entorno de la plataforma \parencite{androidPruebasLocales}. AndroidX Test y Compose
UI Test proporcionan el soporte de las pruebas instrumentadas y de interfaz.
Android Lint complementa esas ejecuciones con análisis estático.

Los scripts PowerShell reúnen tareas de compilación y verificación y conservan
reportes, paquetes y hashes para identificar lo producido. Otros scripts Python
y JavaScript apoyan la preparación de recursos de audio e imagen y las
comprobaciones de notación. Su función es hacer ejecutables operaciones
recurrentes y conservar sus salidas. Los casos, condiciones y resultados de
cada verificación corresponden al capítulo~\ref{cap:pruebas}.

Para producir ilustraciones estáticas también se utilizó una herramienta integrada
de generación de imágenes. El inventario de recursos registra originales,
prompts y exportaciones, y los scripts Python preparan las copias para Android.
Esa procedencia se conserva por recurso, junto con las licencias de fuentes y
bibliotecas y la identificación de materiales aportados por mí.

\section{Organización, relevamiento y escritura}
\label{sec:herramientas-documentacion}
Utilicé Git para el historial de cambios y GitHub para alojar el repositorio,
organizar tareas y especificaciones mediante Issues y conservar entregas.
La interfaz de línea de comandos \texttt{gh} permite consultar y actualizar
esas tareas desde el entorno de trabajo. Los documentos Markdown reúnen
definiciones, decisiones y registros junto al código. Estos recursos sostienen
la trazabilidad explicada en la sección~\ref{sec:metodo-desarrollo}.

Para el relevamiento preparé la encuesta y el formulario de contacto independiente
en Google Forms. El repositorio conserva un generador Apps Script como apoyo a
su preparación, además del cuestionario y las instrucciones de revisión. El propósito,
la población y el análisis previsto se describen en la
sección~\ref{sec:metodo-encuesta}.

Para redactar el informe adopté LaTeX por capítulos. La compilación utiliza
XeLaTeX de la distribución \mbox{MiKTeX}. LaTeX proporciona composición de documentos
técnicos y científicos
\parencite{latexIntroduccion}; en el proyecto permite mantener una fuente editable
con referencias cruzadas y salida PDF. BibLaTeX y Biber gestionan las citas y
referencias con el estilo APA configurado para el informe. Un script PowerShell
automatiza la compilación, mientras que Python y \texttt{pdftotext} permiten
extraer las guías institucionales con referencias a sus páginas. Markdown se
reserva para planificación, procedencia de fuentes y seguimiento de escritura.

\section{Asistentes y agentes de inteligencia artificial}
\label{sec:herramientas-ia}
Al comienzo del proyecto utilicé Gemini Deep Research desde Gemini para explorar
conceptos, localizar antecedentes y obtener informes de investigación.
La herramienta ofrece investigación asistida a partir de consultas y fuentes;
el procedimiento que seguí y la lectura complementaria se detallan en la
sección~\ref{sec:metodo-revision} \parencite{geminiDeepResearch}.

Comencé el desarrollo con Claude Code CLI y posteriormente migré a Codex CLI.
Ambas herramientas permiten trabajar desde la terminal sobre el repositorio,
examinar archivos, realizar cambios y ejecutar comandos
\parencite{anthropicClaudeCode,openaiCodexCli}.
Las utilicé para implementación y revisión de código, discusión de decisiones
y trabajo de escritura. En la etapa actual combino Codex CLI y GrokBot en
simultáneo.

Trabajé principalmente con los modelos que identifiqué como Opus 5.5 en Claude
Code y GPT 6.1 Sol en Codex, alternando con modelos más livianos para tareas
sencillas. Para tareas complejas, diseño y planificación utilicé configuraciones
de razonamiento alto cuando estaban disponibles. La elección respondía a la
complejidad percibida de cada tarea; no realicé una evaluación comparativa
controlada entre modelos. La revisión de aportes, el contraste entre asistentes
y la responsabilidad sobre su aceptación se describen en el
capítulo~\ref{cap:metodologia}.

\section{Configuración técnica de referencia}
\label{sec:herramientas-versiones}
La tabla~\ref{tab:versiones-herramientas} resume valores declarados en los archivos
de construcción y recursos de Neuma 0.12.20, consultados el 4 de noviembre de 2026.
El número de versión de una dependencia, el de un conjunto BOM y el nivel de API
Android identifican elementos diferentes de la configuración. Las versiones de
intérpretes y aplicaciones auxiliares que no se reconstruyeron se conservan fuera
de esta tabla.

\clearpage
\begingroup
\small
\singlespacing
\setlength{\tabcolsep}{4pt}
\setlength{\LTcapwidth}{\textwidth}
\begin{longtable}{p{0.34\textwidth}p{0.18\textwidth}p{0.38\textwidth}}
\caption{Configuración técnica de referencia de Neuma 0.12.20.}
\label{tab:versiones-herramientas}\\
\toprule
Elemento & Versión o nivel & Función \\
\midrule
\endfirsthead
\toprule
Elemento & Versión o nivel & Función \\
\midrule
\endhead
\bottomrule
\endfoot
Android SDK & 26 / 36 / 36 & API mínima, compilación y destino, respectivamente. \\
Kotlin y plugins Kotlin & 2.2.10 & Lenguaje y soporte de Compose y serialización. \\
Compose BOM & 2025.08.01 & Conjunto de versiones de bibliotecas Compose. \\
Lifecycle & 2.9.2 & Estado y ciclo de vida. \\
Corrutinas Kotlin & 1.10.2 & Operaciones asíncronas y pruebas. \\
Serialization JSON & 1.9.0 & Conversión de datos al formato JSON. \\
VexFlow & 5.0.0 & Grabado musical en Rhythm Tap. \\
Gradle Wrapper & 8.14.3 & Distribución de Gradle fijada en el repositorio. \\
Android Gradle Plugin & 8.11.0 & Construcción y empaquetado. \\
Destino de bytecode JVM & 17 & Compatibilidad Java/Kotlin. \\
JUnit & 4.13.2 & Pruebas locales. \\
AndroidX JUnit / runner & 1.2.1 / 1.6.2 & Soporte de pruebas instrumentadas. \\
\end{longtable}
\endgroup
